%============================================================
%  Bd. 14 - Rationale Zahlen
%============================================================

\documentclass[main.tex]{subfiles}

\ifSubfilesClassLoaded{
  \BandSetupStandalone{B14}
}{
}

\title{Bd. 14 - Rationale Zahlen}
\author{Martin Kunze}
\date{}
\setcounter{file}{14}

\hypersetup{
  pdftitle={Bd. 14 - Rationale Zahlen},
  pdfauthor={Martin Kunze}
}

% Die Notationen bleiben bis zu ihrer Aufnahme in die gemeinsame
% Makrodatei lokal in diesem Band.
\providecommand{\PositiveIntegers}{\ensuremath{\mathbb Z_{>0}}}
\providecommand{\RatPairCarrier}{\ensuremath{D_{\mathbb Q}}}
\providecommand{\RatPairRel}{\ensuremath{\mathrel{\sim_{\mathbb Q}}}}
\providecommand{\RatClass}[2]{\ensuremath{\left[#1,#2\right]_{\mathbb Q}}}
\providecommand{\RatEmbed}{\ensuremath{\iota_{\mathbb Q}}}
\providecommand{\RatOrderGraph}{\ensuremath{R^{\leq}_{\mathbb Q}}}
\providecommand{\RatRecipRep}[2]{%
  \ensuremath{\operatorname{rec}_{\mathbb Q}\!\left(#1,#2\right)}%
}
\providecommand{\RationalOrderedField}[1]{%
  \ensuremath{\mathsf{ArchOF}\!\left(#1\right)}%
}

\begin{document}

\title{Bd. 14 - Rationale Zahlen}
\author{Martin Kunze}
\date{}
\hypersetup{pageanchor=false}
\maketitle
\hypersetup{pageanchor=true}
\setcounter{tocdepth}{3}
\tableofcontents
\setcounter{file}{14}
\renewcommand{\FormulaBandID}{14}

\chapter{Einleitung}

Die rationalen Zahlen entstehen aus ganzzahligen Zählern und positiven
ganzzahligen Nennern. Ein Paar \((a,b)\) wird als formaler Bruch
\(a/b\) gelesen. Verschiedene Bruchpaare beschreiben dieselbe Zahl; daher
werden genau diejenigen Paare identifiziert, deren Kreuzprodukte
übereinstimmen. Die Äquivalenz- und Quotiententheorie aus Band~10 sowie die
ganzen Zahlen aus Band~13 liefern die dafür benötigten Grundlagen.

Wie bei der Konstruktion der ganzen Zahlen wird jede Operation erst nach dem
Nachweis ihrer Repräsentantenunabhängigkeit eingeführt. Positive Nenner
vermeiden ein zusätzliches Vorzeichen im Nenner und liefern zugleich eine
einfache Quotientenordnung. Abschließend werden die gewonnenen Gesetze als
archimedisch geordnete Körperstruktur zusammengefasst. Diese Zusammenfassung
ist eine metasprachliche Bündelung der zuvor einzeln bewiesenen Aussagen; sie
setzt keinen bereits eingeführten abstrakten Körperbegriff voraus.

\chapter{Konstruktion der rationalen Zahlen}

\section{Positive ganze Zahlen und Kürzungslemmata}

\FormulaDefDeltaK[Positive ganze Zahlen]{%
  \PositiveIntegers\coloneqq
  \{b\in\mathbb Z\mid 0_{\mathbb Z}<b\}%
}{PositiveIntegersDef}{%
  \DeltaRow{Mengen}{\PositiveIntegers}%
  \DeltaRow{Relationsoperatoren}{<}%
  \DeltaRow{Neue Symbole}{\PositiveIntegers}%
}

\FormulaThmDeltaK[Die Ganzzahleins ist positiv]{%
  1_{\mathbb Z}\in\PositiveIntegers%
}{IntegerOnePositive}{%
  \DeltaRow{Mengen}{\PositiveIntegers}%
}

\emph{Beweisskizze.}
Nach \FormulaRefAuto{IntegerOneDef}[def] ist
\(1_{\mathbb Z}=\IntEmbed(1)\). Die strikte Positivität folgt aus der
Nachfolgerorientierung der Ganzzahlordnung und der Verträglichkeit der
natürlichen Einbettung mit dem Nachfolger.

\FormulaThmDeltaK[Positive ganze Zahlen sind von null verschieden]{%
  b\in\PositiveIntegers\vdash b\neq0_{\mathbb Z}%
}{PositiveIntegerNonzero}{%
  \DeltaRow{Mengen}{b}%
}

\emph{Beweisskizze.}
Aus \(0_{\mathbb Z}<b\) und \(b=0_{\mathbb Z}\) entstünde unmittelbar
\(0_{\mathbb Z}<0_{\mathbb Z}\), im Widerspruch zur Irreflexivität der
strikten Ordnung.

\FormulaThmDeltaK[Die Ganzzahleins ist kleinste positive ganze Zahl]{%
  b\in\PositiveIntegers\vdash1_{\mathbb Z}\leq b%
}{PositiveIntegerOneLowerBound}{%
  \DeltaRow{Mengen}{b}%
  \DeltaRow{Relationsoperatoren}{\leq}%
}

\emph{Beweisskizze.}
Die Normalform \FormulaRefAuto{IntegerSignedNormalForm}[thm] stellt ein
positives \(b\) als \(\IntEmbed(n)\) dar. Wegen \(b\neq0_{\mathbb Z}\) ist
\(n\neq0\), also ein Nachfolger. Die natürliche Ordnung und ihre in Band~13
konstruierte Ganzzahleinbettung liefern daher
\(1_{\mathbb Z}\leq\IntEmbed(n)=b\).

\FormulaThmDeltaK[Abgeschlossenheit positiver ganzer Zahlen unter Multiplikation]{%
  b\in\PositiveIntegers\dsep d\in\PositiveIntegers
  \vdash b\cdot d\in\PositiveIntegers%
}{PositiveIntegerMulClosure}{%
  \DeltaRow{Mengen}{b\dsep d}%
}

\FormulaThmDeltaK[Kürzung eines positiven ganzzahligen Faktors]{%
  a,c\in\mathbb Z\dsep d\in\PositiveIntegers\dsep
  a\cdot d=c\cdot d\vdash a=c%
}{IntegerPositiveFactorCancellation}{%
  \DeltaRow{Mengen}{a\dsep c\dsep d}%
}

\FormulaThmDeltaK[Positive ganzzahlige Faktoren bewahren und reflektieren die Ordnung]{%
  a,c\in\mathbb Z\dsep d\in\PositiveIntegers
  \vdash
  \bigl(a\leq c\leftrightarrow a\cdot d\leq c\cdot d\bigr)%
}{IntegerPositiveFactorOrder}{%
  \DeltaRow{Mengen}{a\dsep c\dsep d}%
}

\emph{Gemeinsame Beweisskizze.}
Nach der Ganzzahlnormalform
\FormulaRefAuto{IntegerSignedNormalForm}[thm] wird jeder Faktor als
eingebettete natürliche Zahl oder als deren Negatives dargestellt. Auf dem
positiven Strahl folgen Produktspositivität und strikte Monotonie durch
Induktion aus der rekursiven Multiplikation auf \(\mathbb N\), der
Distributivität und der strikten Monotonie der Addition. Die multiplikative
Verträglichkeit und Injektivität der Einbettung übertragen diese Aussagen auf
\(\mathbb Z\). Die Kürzungsregel ist der Gleichheitsfall der reflektierten
Ordnung.

\section{Bruchpaare}

\FormulaDefDeltaK[Träger der Bruchpaare]{%
  \RatPairCarrier\coloneqq\mathbb Z\times\PositiveIntegers%
}{RationalPairCarrier}{%
  \DeltaRow{Mengen}{\RatPairCarrier\dsep\PositiveIntegers}%
  \DeltaRow{Neue Symbole}{\RatPairCarrier}%
}

\FormulaDefDeltaK[Äquivalenz von Bruchpaaren]{%
  a,c\in\mathbb Z\dsep b,d\in\PositiveIntegers\vdash
  (a,b)\RatPairRel(c,d)
  \coloneqq a\cdot d=c\cdot b%
}{RationalPairRelDef}{%
  \DeltaRow{Mengen}{a\dsep b\dsep c\dsep d}%
  \DeltaRow{Relationsoperatoren}{\RatPairRel}%
  \DeltaRow{Trägermenge}{\RatPairCarrier}%
  \DeltaRow{Neue Symbole}{\RatPairRel}%
}

\FormulaThmDeltaK[Charakterisierung der Bruchpaaräquivalenz]{%
  a,c\in\mathbb Z\dsep b,d\in\PositiveIntegers\vdash
  \bigl((a,b)\RatPairRel(c,d)
  \leftrightarrow a\cdot d=c\cdot b\bigr)%
}{RationalPairRelChar}{%
  \DeltaRow{Mengen}{a\dsep b\dsep c\dsep d}%
  \DeltaRow{Relationsoperatoren}{\RatPairRel}%
}
\begin{tabproof}
  \proofstep{1}{a,c\in\mathbb Z}{\rA}
  \proofstep{2}{b,d\in\PositiveIntegers}{\rA}
  \proofstep{1,2}{%
    (a,b)\RatPairRel(c,d)\leftrightarrow a\cdot d=c\cdot b}{%
    \FormulaRefAuto{RationalPairRelDef}[def]{1,2}}
\end{tabproof}

\FormulaThmDeltaK[Reflexivität auf Bruchpaaren]{%
  a\in\mathbb Z\dsep b\in\PositiveIntegers
  \vdash(a,b)\RatPairRel(a,b)%
}{RationalPairRelReflexive}{%
  \DeltaRow{Mengen}{a\dsep b}%
  \DeltaRow{Relationsoperatoren}{\RatPairRel}%
}
\begin{tabproof}
  \proofstep{1}{a\in\mathbb Z}{\rA}
  \proofstep{2}{b\in\PositiveIntegers}{\rA}
  \proofstep{}{a\cdot b=a\cdot b}{\rII}
  \proofstep{1,2}{(a,b)\RatPairRel(a,b)}{%
    \FormulaRefAuto{RationalPairRelChar}[thm]{1,2,3}}
\end{tabproof}

\FormulaThmDeltaK[Symmetrie auf Bruchpaaren]{%
  a,c\in\mathbb Z\dsep b,d\in\PositiveIntegers\dsep
  (a,b)\RatPairRel(c,d)
  \vdash(c,d)\RatPairRel(a,b)%
}{RationalPairRelSymmetric}{%
  \DeltaRow{Mengen}{a\dsep b\dsep c\dsep d}%
  \DeltaRow{Relationsoperatoren}{\RatPairRel}%
}
\begin{tabproof}
  \proofstep{1}{a,c\in\mathbb Z}{\rA}
  \proofstep{2}{b,d\in\PositiveIntegers}{\rA}
  \proofstep{3}{(a,b)\RatPairRel(c,d)}{\rA}
  \proofstep{1,2,3}{a\cdot d=c\cdot b}{%
    \FormulaRefAuto{RationalPairRelChar}[thm]{1,2,3}}
  \proofstep{1,2,3}{c\cdot b=a\cdot d}{%
    \FormulaRefAuto{a=b\vdash b=a}{4}}
  \proofstep{1,2,3}{(c,d)\RatPairRel(a,b)}{%
    \FormulaRefAuto{RationalPairRelChar}[thm]{1,2,5}}
\end{tabproof}

\FormulaThmDeltaK[Transitivität auf Bruchpaaren]{%
  \begin{aligned}[t]
    &a,c,e\in\mathbb Z\dsep b,d,f\in\PositiveIntegers\dsep
      (a,b)\RatPairRel(c,d)\dsep{}\\
    &(c,d)\RatPairRel(e,f)
      \vdash(a,b)\RatPairRel(e,f)
  \end{aligned}%
}{RationalPairRelTransitive}{%
  \DeltaRow{Mengen}{a\dsep b\dsep c\dsep d\dsep e\dsep f}%
  \DeltaRow{Relationsoperatoren}{\RatPairRel}%
}

\emph{Beweisskizze.}
Aus \(ad=cb\) und \(cf=ed\) folgt nach Multiplikation mit \(f\)
beziehungsweise \(b\), Kommutativität und Assoziativität
\(d(af)=d(eb)\). Da \(d>0\) ist, liefert
\FormulaRefAuto{IntegerPositiveFactorCancellation}[thm] die Gleichheit
\(af=eb\), also \((a,b)\RatPairRel(e,f)\).

\FormulaThmDeltaK[Die Bruchpaarrelation ist eine Äquivalenzrelation]{%
  \EqRel{\RatPairCarrier,\RatPairRel}%
}{RationalPairEquivalence}{%
  \DeltaRow{Mengen}{\RatPairCarrier}%
  \DeltaRow{Relationsoperatoren}{\RatPairRel}%
}

\emph{Beweisskizze.}
Jedes Element von \(\RatPairCarrier\) besitzt eindeutig die Gestalt
\((a,b)\) mit \(a\in\mathbb Z\) und \(b\in\PositiveIntegers\). Damit sind
Reflexivität, Symmetrie und Transitivität genau die drei vorangehenden Sätze.

\section{Quotient und Bruchklassen}

\FormulaDefDeltaK[Menge der rationalen Zahlen]{%
  \mathbb Q\coloneqq
  \QuotientSet{\RatPairCarrier}{\RatPairRel}%
}{RationalsAsQuotient}{%
  \DeltaRow{Mengen}{\RatPairCarrier\dsep\mathbb Q}%
  \DeltaRow{Relationsoperatoren}{\RatPairRel}%
  \DeltaRow{Neue Symbole}{\mathbb Q}%
}

\FormulaDefDeltaK[Klasse eines Bruchpaares]{%
  a\in\mathbb Z\dsep b\in\PositiveIntegers\vdash
  \RatClass{a}{b}\coloneqq
  \EqClass{(a,b)}{\RatPairRel}{\RatPairCarrier}%
}{RationalClassDef}{%
  \DeltaRow{Mengen}{a\dsep b\dsep\RatPairCarrier}%
  \DeltaRow{Relationsoperatoren}{\RatPairRel}%
  \DeltaRow{Neue Symbole}{\RatClass{a}{b}}%
}

\FormulaThmDeltaK[Bruchklassen sind rationale Zahlen]{%
  a\in\mathbb Z\dsep b\in\PositiveIntegers
  \vdash\RatClass{a}{b}\in\mathbb Q%
}{RationalClassInQ}{%
  \DeltaRow{Mengen}{a\dsep b}%
}

\emph{Beweisskizze.}
Das Paar \((a,b)\) liegt in \(\RatPairCarrier\). Mit
\FormulaRefAuto{RationalPairEquivalence}[thm] liegt seine Äquivalenzklasse in
der Quotientenmenge; anschließend werden
\FormulaRefAuto{RationalClassDef}[def] und
\FormulaRefAuto{RationalsAsQuotient}[def] eingesetzt.

\FormulaThmDeltaK[Gleichheit von Bruchklassen]{%
  a,c\in\mathbb Z\dsep b,d\in\PositiveIntegers\vdash
  \bigl(\RatClass{a}{b}=\RatClass{c}{d}
  \leftrightarrow a\cdot d=c\cdot b\bigr)%
}{RationalClassEquality}{%
  \DeltaRow{Mengen}{a\dsep b\dsep c\dsep d}%
  \DeltaRow{Relationsoperatoren}{\RatPairRel}%
}

\emph{Beweisskizze.}
Für Äquivalenzrelationen sind zwei Klassen genau dann gleich, wenn ihre
Repräsentanten äquivalent sind. Dies folgt aus
\FormulaRefAuto{EqClassEqualIffRel}[thm]. Die Behauptung ergibt sich danach
aus \FormulaRefAuto{RationalPairRelChar}[thm].

\FormulaThmDeltaK[Jede rationale Zahl besitzt ein Bruchpaar]{%
  q\in\mathbb Q\vdash
  \exists a\in\mathbb Z\,\exists b\in\PositiveIntegers\,
  q=\RatClass{a}{b}%
}{RationalPairRepresentative}{%
  \DeltaRow{Mengen}{q\dsep a\dsep b\dsep\RatPairCarrier}%
}

\emph{Beweisskizze.}
Nach der Repräsentanteneigenschaft einer Quotientenmenge besitzt \(q\) eine
Darstellung als Klasse eines Elements von \(\RatPairCarrier\). Das
kartesische Produkt zerlegt dieses Element in einen ganzzahligen Zähler und
einen positiven ganzzahligen Nenner.

\section{Einbettung der ganzen Zahlen}

\FormulaDefDeltaK[Funktionsvorschrift der ganzzahligen Einbettung]{%
  z\in\mathbb Z\vdash
  \RatEmbed(z)\coloneqq\RatClass{z}{1_{\mathbb Z}}%
}{IntegerRationalEmbedding}{%
  \DeltaRow{Mengen}{z}%
  \DeltaRow{Funktionensymbol}{\RatEmbed}%
  \DeltaRow{Neue Symbole}{\RatEmbed}%
}

\FormulaThmDeltaK[Typisierung der ganzzahligen Einbettung]{%
  \RatEmbed\colon\mathbb Z\to\mathbb Q%
}{IntegerRationalEmbeddingFunction}{%
  \DeltaRow{Funktion}{\RatEmbed}%
}

\emph{Beweisskizze.}
Für jedes \(z\in\mathbb Z\) ist \(1_{\mathbb Z}\) positiv und daher
\(\RatClass{z}{1_{\mathbb Z}}\in\mathbb Q\). Das allgemeine
Funktionsbildungstheorem aus Band~05 liefert die eindeutig bestimmte
Abbildung.

\FormulaThmDeltaK[Injektivität der ganzzahligen Einbettung]{%
  \RatEmbed\colon\mathbb Z\inj\mathbb Q%
}{IntegerRationalEmbeddingInjective}{%
  \DeltaRow{Mengen}{z\dsep w}%
  \DeltaRow{Funktion}{\RatEmbed}%
}

\emph{Beweisskizze.}
Aus \(\RatEmbed(z)=\RatEmbed(w)\) folgt nach dem Klassenkriterium
\(z1_{\mathbb Z}=w1_{\mathbb Z}\), also \(z=w\). Zusammen mit der
Typisierung ist dies die Injektivitätsbedingung.

\FormulaDefDeltaK[Null der rationalen Zahlen]{%
  0_{\mathbb Q}\coloneqq\RatEmbed(0_{\mathbb Z})%
}{RationalZeroDef}{%
  \DeltaRow{Neue Symbole}{0_{\mathbb Q}}%
  \DeltaRow{Funktionensymbol}{\RatEmbed}%
}

\FormulaDefDeltaK[Eins der rationalen Zahlen]{%
  1_{\mathbb Q}\coloneqq\RatEmbed(1_{\mathbb Z})%
}{RationalOneDef}{%
  \DeltaRow{Neue Symbole}{1_{\mathbb Q}}%
  \DeltaRow{Funktionensymbol}{\RatEmbed}%
}

\FormulaThmDeltaK[Null- und Einsklasse]{%
  0_{\mathbb Q}=\RatClass{0_{\mathbb Z}}{1_{\mathbb Z}}
  \dsep
  1_{\mathbb Q}=\RatClass{1_{\mathbb Z}}{1_{\mathbb Z}}%
}{RationalZeroOneClasses}{%
  \DeltaRow{Mengen}{\mathbb Q}%
}

\FormulaThmDeltaK[Null und Eins sind rationale Zahlen]{%
  0_{\mathbb Q}\in\mathbb Q\dsep1_{\mathbb Q}\in\mathbb Q%
}{RationalZeroOneCarrier}{%
  \DeltaRow{Mengen}{\mathbb Q}%
}

\emph{Beweisskizze.}
Beide Aussagen folgen unmittelbar aus den beiden Definitionen, der
Funktionsgleichung von \(\RatEmbed\) und
\FormulaRefAuto{IntegerRationalEmbeddingFunction}[thm].

\chapter{Arithmetik der rationalen Zahlen}

\section{Negation}

\FormulaThmDeltaK[Repräsentantenunabhängigkeit der rationalen Negation]{%
  a,c\in\mathbb Z\dsep b,d\in\PositiveIntegers\dsep
  \RatClass{a}{b}=\RatClass{c}{d}
  \vdash\RatClass{-a}{b}=\RatClass{-c}{d}%
}{RationalNegationCompatible}{%
  \DeltaRow{Mengen}{a\dsep b\dsep c\dsep d}%
}

\emph{Beweisskizze.}
Aus \(ad=cb\) folgt durch Negation beider Seiten
\((-a)d=(-c)b\). Das Klassenkriterium liefert die Behauptung.

\FormulaThmDeltaK[Quotientenabstieg der rationalen Negation]{%
  q\in\mathbb Q\vdash
  \exists!u\in\mathbb Q\,
  \exists a\in\mathbb Z\,\exists b\in\PositiveIntegers\,
  \bigl(q=\RatClass{a}{b}\land u=\RatClass{-a}{b}\bigr)%
}{RationalNegationQuotientDescent}{%
  \DeltaRow{Mengen}{q\dsep u\dsep a\dsep b}%
}

\FormulaDefDeltaK[Negation rationaler Zahlen]{%
  q\in\mathbb Q\vdash
  -q\coloneqq\iota u\Bigl(
    u\in\mathbb Q\land
    \exists a\in\mathbb Z\,\exists b\in\PositiveIntegers\,
    (q=\RatClass{a}{b}\land u=\RatClass{-a}{b})
  \Bigr)%
}{RationalNegationDef}{%
  \DeltaRow{Mengen}{q\dsep u\dsep a\dsep b}%
  \DeltaRow{Neue Symbole}{-q}%
}

\FormulaThmDeltaK[Negation auf Bruchklassen]{%
  a\in\mathbb Z\dsep b\in\PositiveIntegers\vdash
  -\RatClass{a}{b}=\RatClass{-a}{b}%
}{RationalNegationClassValue}{%
  \DeltaRow{Mengen}{a\dsep b}%
}

\FormulaThmDeltaK[Abgeschlossenheit und doppelte Negation]{%
  q\in\mathbb Q\vdash
  -q\in\mathbb Q\dsep-(-q)=q%
}{RationalNegationLaws}{%
  \DeltaRow{Mengen}{q}%
}

\emph{Beweisskizze.}
Existenz folgt aus dem Repräsentantensatz, Eindeutigkeit aus der
Repräsentantenunabhängigkeit. Auf einer Bruchklasse werden beide Aussagen zu
\(-(-a)=a\) in \(\mathbb Z\).

\section{Addition}

\FormulaThmDeltaK[Repräsentantenunabhängigkeit der rationalen Addition]{%
  \begin{aligned}[t]
    &a,a',c,c'\in\mathbb Z\dsep
      b,b',d,d'\in\PositiveIntegers\dsep
      \RatClass{a}{b}=\RatClass{a'}{b'}\dsep{}\\
    &\RatClass{c}{d}=\RatClass{c'}{d'}\vdash
      \RatClass{a\cdot d+c\cdot b}{b\cdot d}
      =\RatClass{a'\cdot d'+c'\cdot b'}{b'\cdot d'}
  \end{aligned}%
}{RationalAdditionCompatible}{%
  \DeltaRow{Mengen}{a\dsep a'\dsep b\dsep b'\dsep
    c\dsep c'\dsep d\dsep d'}%
}

\emph{Beweisskizze.}
Die Voraussetzungen sind \(ab'=a'b\) und \(cd'=c'd\). Nach
Ausmultiplizieren der beiden Kreuzprodukte werden genau diese beiden
Gleichheiten eingesetzt; Assoziativität, Kommutativität und Distributivität
in \(\mathbb Z\) ordnen die Faktoren. Die Nennerprodukte sind nach
\FormulaRefAuto{PositiveIntegerMulClosure}[thm] positiv.

\FormulaThmDeltaK[Quotientenabstieg der rationalen Addition]{%
  q,r\in\mathbb Q\vdash
  \exists!u\in\mathbb Q\,
  \begin{aligned}[t]
    \exists a,c\in\mathbb Z\,\exists b,d\in\PositiveIntegers\,
    \bigl(&q=\RatClass{a}{b}\land r=\RatClass{c}{d}\land{}\\
          &u=\RatClass{a\cdot d+c\cdot b}{b\cdot d}\bigr)
  \end{aligned}%
}{RationalAdditionQuotientDescent}{%
  \DeltaRow{Mengen}{q\dsep r\dsep u\dsep a\dsep b\dsep c\dsep d}%
}

\FormulaDefDeltaK[Addition rationaler Zahlen]{%
  q,r\in\mathbb Q\vdash
  q+r\coloneqq\iota u\Bigl(
  \begin{aligned}[t]
    &u\in\mathbb Q\land
      \exists a,c\in\mathbb Z\,\exists b,d\in\PositiveIntegers\,\bigl({}\\
    &q=\RatClass{a}{b}\land r=\RatClass{c}{d}\land
      u=\RatClass{a\cdot d+c\cdot b}{b\cdot d}\bigr)
  \end{aligned}
  \Bigr)%
}{RationalAdditionDef}{%
  \DeltaRow{Mengen}{q\dsep r\dsep u\dsep a\dsep b\dsep c\dsep d}%
  \DeltaRow{Neue Symbole}{+}%
}

\FormulaThmDeltaK[Addition auf Bruchklassen]{%
  a,c\in\mathbb Z\dsep b,d\in\PositiveIntegers\vdash
  \RatClass{a}{b}+\RatClass{c}{d}
  =\RatClass{a\cdot d+c\cdot b}{b\cdot d}%
}{RationalAdditionClassValue}{%
  \DeltaRow{Mengen}{a\dsep b\dsep c\dsep d}%
}

\FormulaThmDeltaK[Abgeschlossenheit der rationalen Addition]{%
  q,r\in\mathbb Q\vdash q+r\in\mathbb Q%
}{RationalAddClosure}{%
  \DeltaRow{Mengen}{q\dsep r}%
}

\section{Multiplikation}

\FormulaThmDeltaK[Repräsentantenunabhängigkeit der rationalen Multiplikation]{%
  \begin{aligned}[t]
    &a,a',c,c'\in\mathbb Z\dsep
      b,b',d,d'\in\PositiveIntegers\dsep
      \RatClass{a}{b}=\RatClass{a'}{b'}\dsep{}\\
    &\RatClass{c}{d}=\RatClass{c'}{d'}\vdash
      \RatClass{a\cdot c}{b\cdot d}
      =\RatClass{a'\cdot c'}{b'\cdot d'}
  \end{aligned}%
}{RationalMultiplicationCompatible}{%
  \DeltaRow{Mengen}{a\dsep a'\dsep b\dsep b'\dsep
    c\dsep c'\dsep d\dsep d'}%
}

\emph{Beweisskizze.}
Aus \(ab'=a'b\) und \(cd'=c'd\) folgt durch Multiplikation und Umordnung
\((ac)(b'd')=(a'c')(bd)\). Die Nennerprodukte sind positiv, sodass das
Klassenkriterium anwendbar ist.

\FormulaThmDeltaK[Quotientenabstieg der rationalen Multiplikation]{%
  q,r\in\mathbb Q\vdash
  \exists!u\in\mathbb Q\,
  \begin{aligned}[t]
    \exists a,c\in\mathbb Z\,\exists b,d\in\PositiveIntegers\,
    \bigl(&q=\RatClass{a}{b}\land r=\RatClass{c}{d}\land{}\\
          &u=\RatClass{a\cdot c}{b\cdot d}\bigr)
  \end{aligned}%
}{RationalMultiplicationQuotientDescent}{%
  \DeltaRow{Mengen}{q\dsep r\dsep u\dsep a\dsep b\dsep c\dsep d}%
}

\FormulaDefDeltaK[Multiplikation rationaler Zahlen]{%
  q,r\in\mathbb Q\vdash
  q\cdot r\coloneqq\iota u\Bigl(
  \begin{aligned}[t]
    &u\in\mathbb Q\land
      \exists a,c\in\mathbb Z\,\exists b,d\in\PositiveIntegers\,\bigl({}\\
    &q=\RatClass{a}{b}\land r=\RatClass{c}{d}\land
      u=\RatClass{a\cdot c}{b\cdot d}\bigr)
  \end{aligned}
  \Bigr)%
}{RationalMultiplicationDef}{%
  \DeltaRow{Mengen}{q\dsep r\dsep u\dsep a\dsep b\dsep c\dsep d}%
  \DeltaRow{Neue Symbole}{\cdot}%
}

\FormulaThmDeltaK[Multiplikation auf Bruchklassen]{%
  a,c\in\mathbb Z\dsep b,d\in\PositiveIntegers\vdash
  \RatClass{a}{b}\cdot\RatClass{c}{d}
  =\RatClass{a\cdot c}{b\cdot d}%
}{RationalMultiplicationClassValue}{%
  \DeltaRow{Mengen}{a\dsep b\dsep c\dsep d}%
}

\FormulaThmDeltaK[Abgeschlossenheit der rationalen Multiplikation]{%
  q,r\in\mathbb Q\vdash q\cdot r\in\mathbb Q%
}{RationalMulClosure}{%
  \DeltaRow{Mengen}{q\dsep r}%
}

\section{Grundlegende Rechengesetze}

\FormulaThmDeltaK[Assoziativität der rationalen Addition]{%
  p,q,r\in\mathbb Q\vdash(p+q)+r=p+(q+r)%
}{RationalAddAssociative}{%
  \DeltaRow{Mengen}{p\dsep q\dsep r}%
}

\FormulaThmDeltaK[Kommutativität der rationalen Addition]{%
  q,r\in\mathbb Q\vdash q+r=r+q%
}{RationalAddCommutative}{%
  \DeltaRow{Mengen}{q\dsep r}%
}

\FormulaThmDeltaK[Nullgesetz der rationalen Addition]{%
  q\in\mathbb Q\vdash q+0_{\mathbb Q}=q%
}{RationalAddZero}{%
  \DeltaRow{Mengen}{q}%
}

\FormulaThmDeltaK[Additives Inverses rationaler Zahlen]{%
  q\in\mathbb Q\vdash q+(-q)=0_{\mathbb Q}%
}{RationalAddInverse}{%
  \DeltaRow{Mengen}{q}%
}

\FormulaThmDeltaK[Assoziativität der rationalen Multiplikation]{%
  p,q,r\in\mathbb Q\vdash(p\cdot q)\cdot r=p\cdot(q\cdot r)%
}{RationalMulAssociative}{%
  \DeltaRow{Mengen}{p\dsep q\dsep r}%
}

\FormulaThmDeltaK[Kommutativität der rationalen Multiplikation]{%
  q,r\in\mathbb Q\vdash q\cdot r=r\cdot q%
}{RationalMulCommutative}{%
  \DeltaRow{Mengen}{q\dsep r}%
}

\FormulaThmDeltaK[Einheitsgesetz der rationalen Multiplikation]{%
  q\in\mathbb Q\vdash q\cdot1_{\mathbb Q}=q%
}{RationalMulOne}{%
  \DeltaRow{Mengen}{q}%
}

\FormulaThmDeltaK[Nullgesetz der rationalen Multiplikation]{%
  q\in\mathbb Q\vdash q\cdot0_{\mathbb Q}=0_{\mathbb Q}%
}{RationalMulZero}{%
  \DeltaRow{Mengen}{q}%
}

\FormulaThmDeltaK[Distributivität rationaler Zahlen]{%
  p,q,r\in\mathbb Q\vdash
  p\cdot(q+r)=p\cdot q+p\cdot r%
}{RationalDistributive}{%
  \DeltaRow{Mengen}{p\dsep q\dsep r}%
}

\emph{Gemeinsame Beweisskizze.}
Man wählt Bruchdarstellungen aller beteiligten rationalen Zahlen und setzt die
Klassenformeln für Addition, Multiplikation und Negation ein. Nach
Ausmultiplizieren reduzieren sich die Aussagen auf die entsprechenden
Ganzzahlgesetze
\FormulaRefAuto{IntAddAssociative}[thm],
\FormulaRefAuto{IntAddCommutative}[thm],
\FormulaRefAuto{IntMulAssociative}[thm],
\FormulaRefAuto{IntMulCommutative}[thm] und
\FormulaRefAuto{IntMulLeftDistributive}[thm]. Die Null- und Einsgesetze folgen
aus den eingebetteten Klassen mit Nenner \(1_{\mathbb Z}\).

\FormulaThmDeltaK[Additivität der ganzzahligen Einbettung]{%
  z,w\in\mathbb Z\vdash
  \RatEmbed(z+w)=\RatEmbed(z)+\RatEmbed(w)%
}{IntegerRationalEmbeddingAdditive}{%
  \DeltaRow{Mengen}{z\dsep w}%
  \DeltaRow{Funktion}{\RatEmbed}%
}

\FormulaThmDeltaK[Multiplikativität der ganzzahligen Einbettung]{%
  z,w\in\mathbb Z\vdash
  \RatEmbed(z\cdot w)=\RatEmbed(z)\cdot\RatEmbed(w)%
}{IntegerRationalEmbeddingMultiplicative}{%
  \DeltaRow{Mengen}{z\dsep w}%
  \DeltaRow{Funktion}{\RatEmbed}%
}

\emph{Beweisskizze.}
Beide Seiten werden als Bruchklassen mit Nenner \(1_{\mathbb Z}\)
geschrieben. Die Klassenformeln und die Einheitsgesetze in \(\mathbb Z\)
liefern jeweils dieselbe Klasse.

\section{Kehrwert und Division}

\FormulaThmDeltaK[Nullkriterium einer Bruchklasse]{%
  a\in\mathbb Z\dsep b\in\PositiveIntegers\vdash
  \bigl(\RatClass{a}{b}=0_{\mathbb Q}
  \leftrightarrow a=0_{\mathbb Z}\bigr)%
}{RationalClassZeroIffNumeratorZero}{%
  \DeltaRow{Mengen}{a\dsep b}%
}

\emph{Beweisskizze.}
Mit der Nullklasse lautet das Klassenkriterium
\(a1_{\mathbb Z}=0_{\mathbb Z}b\). Die Einheits- und Nullgesetze in
\(\mathbb Z\) reduzieren dies auf \(a=0_{\mathbb Z}\).

\FormulaDefDeltaK[Reziproker Repräsentantenwert]{%
  \begin{aligned}[t]
    &a\in\mathbb Z\dsep b\in\PositiveIntegers\dsep
      a\neq0_{\mathbb Z}\vdash{}\\
    &\RatRecipRep{a}{b}\coloneqq
      \begin{cases}
        \RatClass{b}{a},&0_{\mathbb Z}<a,\\
        \RatClass{-b}{-a},&a<0_{\mathbb Z}.
      \end{cases}
  \end{aligned}%
}{RationalReciprocalRepresentativeDef}{%
  \DeltaRow{Mengen}{a\dsep b}%
  \DeltaRow{Neue Symbole}{\RatRecipRep{a}{b}}%
}

Wegen der Totalität der Ganzzahlordnung gilt für
\(a\neq0_{\mathbb Z}\) genau einer der beiden Fälle. Im ersten Fall ist
\(a\) selbst ein positiver Nenner, im zweiten Fall ist \(-a\) positiv.

\FormulaThmDeltaK[Typisierung des reziproken Repräsentantenwertes]{%
  a\in\mathbb Z\dsep b\in\PositiveIntegers\dsep
  a\neq0_{\mathbb Z}
  \vdash\RatRecipRep{a}{b}\in\mathbb Q%
}{RationalReciprocalRepresentativeInQ}{%
  \DeltaRow{Mengen}{a\dsep b}%
}

\FormulaThmDeltaK[Repräsentantenunabhängigkeit des Kehrwertes]{%
  \begin{aligned}[t]
    &a,c\in\mathbb Z\dsep b,d\in\PositiveIntegers\dsep
      a\neq0_{\mathbb Z}\dsep c\neq0_{\mathbb Z}\dsep{}\\
    &\RatClass{a}{b}=\RatClass{c}{d}
      \vdash\RatRecipRep{a}{b}=\RatRecipRep{c}{d}
  \end{aligned}%
}{RationalReciprocalCompatible}{%
  \DeltaRow{Mengen}{a\dsep b\dsep c\dsep d}%
}

\emph{Beweisskizze.}
Aus \(ad=cb\) und positiven Nennern folgt, dass \(a\) und \(c\) dasselbe
Vorzeichen besitzen. Sind beide positiv, ist für die reziproken Klassen die
benötigte Kreuzgleichheit \(bc=da\); sie ist nur die symmetrisch umgeordnete
Ausgangsgleichheit. Sind beide negativ, ergibt dieselbe Rechnung nach
Negation aller Zähler und Nenner die Gleichheit der Klassen.

\FormulaThmDeltaK[Quotientenabstieg des Kehrwertes]{%
  \begin{aligned}[t]
    &q\in\mathbb Q\dsep q\neq0_{\mathbb Q}\vdash{}\\
    &\exists!u\in\mathbb Q\,
      \exists a\in\mathbb Z\,\exists b\in\PositiveIntegers\,\bigl({}\\
    &\qquad q=\RatClass{a}{b}\land
      a\neq0_{\mathbb Z}\land u=\RatRecipRep{a}{b}\bigr)
  \end{aligned}%
}{RationalReciprocalQuotientDescent}{%
  \DeltaRow{Mengen}{q\dsep u\dsep a\dsep b}%
}

\FormulaDefDeltaK[Kehrwert einer von null verschiedenen rationalen Zahl]{%
  \begin{aligned}[t]
    &q\in\mathbb Q\dsep q\neq0_{\mathbb Q}\vdash{}\\
    &q^{-1}\coloneqq\iota u\Bigl(
      u\in\mathbb Q\land
      \exists a\in\mathbb Z\,\exists b\in\PositiveIntegers\,\bigl({}\\
    &\qquad q=\RatClass{a}{b}\land a\neq0_{\mathbb Z}\land
      u=\RatRecipRep{a}{b}\bigr)
      \Bigr)
  \end{aligned}%
}{RationalReciprocalDef}{%
  \DeltaRow{Mengen}{q\dsep u\dsep a\dsep b}%
  \DeltaRow{Neue Symbole}{q^{-1}}%
}

\FormulaThmDeltaK[Kehrwert auf Bruchklassen]{%
  a\in\mathbb Z\dsep b\in\PositiveIntegers\dsep
  a\neq0_{\mathbb Z}
  \vdash
  \bigl(\RatClass{a}{b}\bigr)^{-1}=\RatRecipRep{a}{b}%
}{RationalReciprocalClassValue}{%
  \DeltaRow{Mengen}{a\dsep b}%
  \DeltaRow{Funktionsoperatoren}{(\,\cdot\,)^{-1}}%
}

Die Klassenformel folgt unmittelbar aus dem eindeutigen Quotientenabstieg
\FormulaRefAuto{RationalReciprocalQuotientDescent}[thm] und der Definition
\FormulaRefAuto{RationalReciprocalDef}[def].

\FormulaThmDeltaK[Abgeschlossenheit des rationalen Kehrwertes]{%
  q\in\mathbb Q\dsep q\neq0_{\mathbb Q}
  \vdash q^{-1}\in\mathbb Q%
}{RationalReciprocalClosure}{%
  \DeltaRow{Mengen}{q}%
}

\FormulaThmDeltaK[Multiplikatives Inverses rationaler Zahlen]{%
  q\in\mathbb Q\dsep q\neq0_{\mathbb Q}
  \vdash q\cdot q^{-1}=1_{\mathbb Q}%
}{RationalMulInverse}{%
  \DeltaRow{Mengen}{q}%
}

\emph{Beweisskizze.}
Für \(q=\RatClass{a}{b}\) mit \(a>0\) ist das Produkt
\(\RatClass{ab}{ba}\), für \(a<0\) nach zweimaliger Negation dieselbe
Klasse. In beiden Fällen ist sie nach dem Klassenkriterium gleich
\(\RatClass{1_{\mathbb Z}}{1_{\mathbb Z}}\).

\FormulaDefDeltaK[Division rationaler Zahlen]{%
  q,r\in\mathbb Q\dsep r\neq0_{\mathbb Q}\vdash
  \frac qr\coloneqq q\cdot r^{-1}%
}{RationalDivisionDef}{%
  \DeltaRow{Mengen}{q\dsep r}%
  \DeltaRow{Neue Symbole}{\frac qr}%
}

\FormulaThmDeltaK[Abgeschlossenheit der rationalen Division]{%
  q,r\in\mathbb Q\dsep r\neq0_{\mathbb Q}
  \vdash\frac qr\in\mathbb Q%
}{RationalDivisionClosure}{%
  \DeltaRow{Mengen}{q\dsep r}%
}

\chapter{Ordnung der rationalen Zahlen}

\section{Quotientenordnung}

\FormulaThmDeltaK[Repräsentantenunabhängigkeit der rationalen Ordnung]{%
  \begin{aligned}[t]
    &a,a',c,c'\in\mathbb Z\dsep
      b,b',d,d'\in\PositiveIntegers\dsep
      \RatClass{a}{b}=\RatClass{a'}{b'}\dsep{}\\
    &\RatClass{c}{d}=\RatClass{c'}{d'}\vdash
      \bigl(a\cdot d\leq c\cdot b
      \leftrightarrow a'\cdot d'\leq c'\cdot b'\bigr)
  \end{aligned}%
}{RationalOrderCompatible}{%
  \DeltaRow{Mengen}{a\dsep a'\dsep b\dsep b'\dsep
    c\dsep c'\dsep d\dsep d'}%
}

\emph{Beweisskizze.}
Man multipliziert die erste Ungleichung mit dem positiven Produkt
\(b'd'\). Die beiden Klassengleichheiten ersetzen anschließend
\(ab'\) durch \(a'b\) und \(cd'\) durch \(c'd\). Nach Umordnung bleibt die
zweite Ungleichung mit dem gemeinsamen positiven Faktor \(bd\), der mit
\FormulaRefAuto{IntegerPositiveFactorOrder}[thm] reflektiert werden darf.
Die Rückrichtung ist symmetrisch.

\FormulaDefDeltaK[Rationale Ordnung als Quotientenrelation]{%
  \begin{aligned}[t]
  \RatOrderGraph\coloneqq
  \bigl\{&(\RatClass{a}{b},\RatClass{c}{d})
    \in\mathbb Q\times\mathbb Q\bigm|{}\\[-2pt]
    &a,c\in\mathbb Z\land b,d\in\PositiveIntegers
      \land a\cdot d\leq c\cdot b\bigr\}.
  \end{aligned}%
}{RationalOrderRelationDef}{%
  \DeltaRow{Mengen}{a\dsep b\dsep c\dsep d\dsep\RatOrderGraph}%
  \DeltaRow{Neue Symbole}{\RatOrderGraph}%
}

\FormulaThmDeltaK[Quotientenabstieg der rationalen Ordnung]{%
  \begin{aligned}[t]
  &\RatOrderGraph\subseteq\mathbb Q\times\mathbb Q\land{}\\[-2pt]
  &\forall a,c\in\mathbb Z\,\forall b,d\in\PositiveIntegers\,
    \bigl((\RatClass{a}{b},\RatClass{c}{d})\in\RatOrderGraph
      \leftrightarrow a\cdot d\leq c\cdot b\bigr).
  \end{aligned}%
}{RationalOrderQuotientDescent}{%
  \DeltaRow{Mengen}{a\dsep b\dsep c\dsep d\dsep\RatOrderGraph}%
}

Die Existenz der Relation folgt durch Aussonderung in
\(\mathbb Q\times\mathbb Q\). Ihre Klassencharakterisierung ist wegen
\FormulaRefAuto{RationalOrderCompatible}[thm] unabhängig von der Wahl der
Repräsentanten.

\FormulaDefDeltaK[Ordnungsoperator der rationalen Zahlen]{%
  q,r\in\mathbb Q\vdash
  q\leq_{\mathbb Q}r\coloneqq(q,r)\in\RatOrderGraph%
}{RationalOrderDef}{%
  \DeltaRow{Mengen}{q\dsep r\dsep\RatOrderGraph}%
  \DeltaRow{Relationsoperatoren}{\leq_{\mathbb Q}}%
  \DeltaRow{Neue Symbole}{\leq_{\mathbb Q}}%
}

\FormulaDefDeltaK[Notation der rationalen Ordnung]{%
  q,r\in\mathbb Q\vdash
  q\leq r\coloneqq q\leq_{\mathbb Q}r%
}{RationalOrderNotation}{%
  \DeltaRow{Mengen}{q\dsep r}%
  \DeltaRow{Relationsoperatoren}{\leq\dsep\leq_{\mathbb Q}}%
}

\FormulaThmDeltaK[Klassencharakterisierung der rationalen Ordnung]{%
  a,c\in\mathbb Z\dsep b,d\in\PositiveIntegers\vdash
  \bigl(\RatClass{a}{b}\leq\RatClass{c}{d}
  \leftrightarrow a\cdot d\leq c\cdot b\bigr)%
}{RationalOrderClassChar}{%
  \DeltaRow{Mengen}{a\dsep b\dsep c\dsep d}%
  \DeltaRow{Relationsoperatoren}{\leq\dsep\leq_{\mathbb Q}}%
}

\FormulaThmDeltaK[Totale Ordnung der rationalen Zahlen]{%
  \TotOrd{\mathbb Q,\leq_{\mathbb Q}}%
}{RationalOrderTotal}{%
  \DeltaRow{Mengen}{q\dsep r\dsep s}%
  \DeltaRow{Relationsoperatoren}{\leq_{\mathbb Q}}%
}

\emph{Beweisskizze.}
Reflexivität und Totalität folgen unmittelbar aus der totalen
Ganzzahlordnung auf den Kreuzprodukten. Antisymmetrie ist das
Klassenkriterium. Für die Transitivität werden die beiden Ungleichungen mit
den jeweils fehlenden positiven Nennern multipliziert; danach wird der
gemeinsame positive Faktor gekürzt.

\FormulaDefDeltaK[Strikte Ordnung der rationalen Zahlen]{%
  q,r\in\mathbb Q\vdash
  q<r\coloneqq q\leq r\land q\neq r%
}{RationalStrictOrderDef}{%
  \DeltaRow{Mengen}{q\dsep r}%
  \DeltaRow{Relationsoperatoren}{<}%
}

\FormulaThmDeltaK[Strikte Klassencharakterisierung der rationalen Ordnung]{%
  a,c\in\mathbb Z\dsep b,d\in\PositiveIntegers\vdash
  \bigl(\RatClass{a}{b}<\RatClass{c}{d}
  \leftrightarrow a\cdot d<c\cdot b\bigr)%
}{RationalStrictOrderClassChar}{%
  \DeltaRow{Mengen}{a\dsep b\dsep c\dsep d}%
}

\section{Verträglichkeit von Ordnung und Arithmetik}

\FormulaThmDeltaK[Translationsinvarianz der rationalen Ordnung]{%
  p,q,r\in\mathbb Q\vdash
  \bigl(q\leq r\leftrightarrow p+q\leq p+r\bigr)%
}{RationalOrderTranslation}{%
  \DeltaRow{Mengen}{p\dsep q\dsep r}%
}

\FormulaThmDeltaK[Produkt nichtnegativer rationaler Zahlen ist nichtnegativ]{%
  q,r\in\mathbb Q\dsep0_{\mathbb Q}\leq q\dsep0_{\mathbb Q}\leq r
  \vdash0_{\mathbb Q}\leq q\cdot r%
}{RationalNonnegativeProduct}{%
  \DeltaRow{Mengen}{q\dsep r}%
}

\FormulaThmDeltaK[Positive rationale Faktoren bewahren die strikte Ordnung]{%
  p,q,r\in\mathbb Q\dsep q<r\dsep0_{\mathbb Q}<p
  \vdash p\cdot q<p\cdot r%
}{RationalPositiveMulStrictOrder}{%
  \DeltaRow{Mengen}{p\dsep q\dsep r}%
}

\emph{Gemeinsame Beweisskizze.}
Nach Wahl positiver Nenner werden die drei Behauptungen mit der
Klassencharakterisierung in Ganzzahlungleichungen übersetzt. Für die Addition
heben sich die gleichen Kreuzterme auf. Für die Multiplikation werden
Produkte positiver Nenner und die Ganzzahlregeln für positive Faktoren
verwendet.

\FormulaThmDeltaK[Positive rationale Zahlen haben positive Kehrwerte]{%
  q\in\mathbb Q\dsep0_{\mathbb Q}<q
  \vdash0_{\mathbb Q}<q^{-1}%
}{RationalPositiveReciprocal}{%
  \DeltaRow{Mengen}{q\dsep a\dsep b}%
  \DeltaRow{Relationsoperatoren}{<}%
  \DeltaRow{Funktionsoperatoren}{(\,\cdot\,)^{-1}}%
}

\emph{Beweisskizze.}
Sei \(q=\RatClass{a}{b}\). Aus \(0_{\mathbb Q}<q\) folgt mit der strikten
Klassencharakterisierung \(0_{\mathbb Z}<a\). Daher ist nach
\FormulaRefAuto{RationalReciprocalClassValue}[thm]
\(q^{-1}=\RatClass{b}{a}\). Weil auch \(b>0_{\mathbb Z}\) gilt, liefert
dieselbe Klassencharakterisierung \(0_{\mathbb Q}<q^{-1}\).

\FormulaThmDeltaK[Ordnungstreue der ganzzahligen Einbettung]{%
  z,w\in\mathbb Z\vdash
  \bigl(z\leq w\leftrightarrow\RatEmbed(z)\leq\RatEmbed(w)\bigr)%
}{IntegerRationalEmbeddingOrder}{%
  \DeltaRow{Mengen}{z\dsep w}%
  \DeltaRow{Funktion}{\RatEmbed}%
}

\emph{Beweisskizze.}
Beide eingebetteten Zahlen haben Nenner \(1_{\mathbb Z}\). Das
Kreuzproduktkriterium reduziert die rationale Ungleichung deshalb mittels
der Einheitsgesetze auf \(z\leq w\).

\section{Dichte und archimedische Eigenschaft}

\FormulaThmDeltaK[Dichte der rationalen Ordnung]{%
  q,r\in\mathbb Q\dsep q<r
  \vdash\exists s\in\mathbb Q\,(q<s\land s<r)%
}{RationalOrderDense}{%
  \DeltaRow{Mengen}{q\dsep r\dsep s}%
}

\emph{Beweisskizze.}
Die Zahl \(2_{\mathbb Q}\coloneqq1_{\mathbb Q}+1_{\mathbb Q}\) ist positiv
und daher von null verschieden. Für
\(s\coloneqq(q+r)/2_{\mathbb Q}\) folgt aus \(q<r\) durch Addition und
Multiplikation mit dem nach
\FormulaRefAuto{RationalPositiveReciprocal}[thm] positiven Kehrwert von
\(2_{\mathbb Q}\), dass \(q<s<r\).

\FormulaThmDeltaK[Archimedische Eigenschaft der rationalen Zahlen]{%
  q\in\mathbb Q\vdash
  \exists n\in\mathbb N\,
  q<\RatEmbed(\IntEmbed(n))%
}{RationalArchimedean}{%
  \DeltaRow{Mengen}{q\dsep n}%
  \DeltaRow{Funktionensymbole}{\IntEmbed\dsep\RatEmbed}%
}

\emph{Beweisskizze.}
Schreibe \(q=\RatClass{a}{b}\) mit \(b>0\). Nach der Ganzzahlnormalform
ist \(a\) eine eingebettete natürliche Zahl oder deren Negatives. Man wählt
eine positive natürliche Zahl \(n\), deren Einbettung strikt größer als
\(a\) ist. Nach \FormulaRefAuto{PositiveIntegerOneLowerBound}[thm] gilt
\(1_{\mathbb Z}\leq b\); Multiplikation mit dem positiven
\(\IntEmbed(n)\) liefert daher
\(\IntEmbed(n)\leq\IntEmbed(n)b\). Somit ist
\(a<\IntEmbed(n)b\), und das strikte Kreuzproduktkriterium ergibt
\(q<\RatEmbed(\IntEmbed(n))\).

\chapter{Axiomatische Zusammenfassung}

Die folgende Tafel bündelt die für spätere Bände maßgeblichen Resultate.
Der Ausdruck \(\RationalOrderedField{\mathbb Q}\) ist an dieser Stelle nur
ein Name für die aufgeführten Aussagen. Insbesondere werden Dichte und
archimedische Eigenschaft als bewiesene Zusätze registriert und nicht als
Bestandteile der späteren abstrakten Körperdefinition ausgegeben.

\FormulaDefDeltaK[Rationale Zahlen als archimedisch geordneter Körper]{%
  \RationalOrderedField{\mathbb Q}%
}{RationalOrderedFieldAxioms}{%
  \DeltaRow{\makecell[l]{\textbf{Träger und}\\\textbf{Konstanten}}}{}%
  \DeltaRow{Träger}{\mathbb Q}%
  \DeltaRow{Null und Eins}{0_{\mathbb Q},1_{\mathbb Q}\in\mathbb Q}
    [\FormulaRefAuto{RationalZeroOneCarrier}[thm]]%
  \DeltaRow{\makecell[l]{Verschiedene\\Konstanten}}{0_{\mathbb Q}\neq1_{\mathbb Q}}
    [\FormulaRefAuto{IntegerRationalEmbeddingInjective}[thm]]%
  \DeltaRow{\textbf{Addition}}{}%
  \DeltaRow{Abgeschlossenheit}{q,r\in\mathbb Q\vdash q+r\in\mathbb Q}
    [\FormulaRefAuto{RationalAddClosure}[thm]]%
  \DeltaRow{Assoziativität}{p,q,r\in\mathbb Q\vdash(p+q)+r=p+(q+r)}
    [\FormulaRefAuto{RationalAddAssociative}[thm]]%
  \DeltaRow{Kommutativität}{q,r\in\mathbb Q\vdash q+r=r+q}
    [\FormulaRefAuto{RationalAddCommutative}[thm]]%
  \DeltaRow{Negation}{q\in\mathbb Q\vdash-q\in\mathbb Q}
    [\FormulaRefAuto{RationalNegationLaws}[thm]]%
  \DeltaRow{Null}{q\in\mathbb Q\vdash q+0_{\mathbb Q}=q}
    [\FormulaRefAuto{RationalAddZero}[thm]]%
  \DeltaRow{Inverse}{q\in\mathbb Q\vdash q+(-q)=0_{\mathbb Q}}
    [\FormulaRefAuto{RationalAddInverse}[thm]]%
  \DeltaRow{\textbf{Multiplikation}}{}%
  \DeltaRow{Abgeschlossenheit}{q,r\in\mathbb Q\vdash q\cdot r\in\mathbb Q}
    [\FormulaRefAuto{RationalMulClosure}[thm]]%
  \DeltaRow{Assoziativität}{\begin{aligned}[t]
    &p,q,r\in\mathbb Q\vdash{}\\
    &(p\cdot q)\cdot r=p\cdot(q\cdot r)
    \end{aligned}}
    [\FormulaRefAuto{RationalMulAssociative}[thm]]%
  \DeltaRow{Kommutativität}{q,r\in\mathbb Q\vdash q\cdot r=r\cdot q}
    [\FormulaRefAuto{RationalMulCommutative}[thm]]%
  \DeltaRow{Eins}{q\in\mathbb Q\vdash q\cdot1_{\mathbb Q}=q}
    [\FormulaRefAuto{RationalMulOne}[thm]]%
  \DeltaRow{Inverse}{\begin{aligned}[t]
    &q\in\mathbb Q\dsep q\neq0_{\mathbb Q}\vdash{}\\
    &q\cdot q^{-1}=1_{\mathbb Q}
    \end{aligned}}
    [\FormulaRefAuto{RationalMulInverse}[thm]]%
  \DeltaRow{Kehrwert}{\begin{aligned}[t]
    &q\in\mathbb Q\dsep q\neq0_{\mathbb Q}\vdash{}\\
    &q^{-1}\in\mathbb Q
    \end{aligned}}
    [\FormulaRefAuto{RationalReciprocalClosure}[thm]]%
  \DeltaRow{Distributivität}{\begin{aligned}[t]
    &p,q,r\in\mathbb Q\vdash{}\\
    &p\cdot(q+r)=p\cdot q+p\cdot r
    \end{aligned}}
    [\FormulaRefAuto{RationalDistributive}[thm]]%
  \DeltaRow{\textbf{Ordnung}}{}%
  \DeltaRow{Totale Ordnung}{\TotOrd{\mathbb Q,\leq_{\mathbb Q}}}
    [\FormulaRefAuto{RationalOrderTotal}[thm]]%
  \DeltaRow{Translation}{\begin{aligned}[t]
    &p,q,r\in\mathbb Q\vdash{}\\
    &(q\leq r\leftrightarrow p+q\leq p+r)
    \end{aligned}}
    [\FormulaRefAuto{RationalOrderTranslation}[thm]]%
  \DeltaRow{Positive Produkte}{\begin{aligned}[t]
    &q,r\in\mathbb Q\dsep0_{\mathbb Q}\leq q\dsep{}\\
    &0_{\mathbb Q}\leq r\vdash0_{\mathbb Q}\leq q\cdot r
    \end{aligned}}
    [\FormulaRefAuto{RationalNonnegativeProduct}[thm]]%
  \DeltaRow{\textbf{Zusätze}}{}%
  \DeltaRow{Dichte}{\begin{aligned}[t]
    &q,r\in\mathbb Q\dsep q<r\vdash{}\\
    &\exists s\in\mathbb Q\,(q<s\land s<r)
    \end{aligned}}
    [\FormulaRefAuto{RationalOrderDense}[thm]]%
  \DeltaRow{Archimedisch}{\begin{aligned}[t]
    &q\in\mathbb Q\vdash{}\\
    &\exists n\in\mathbb N\,q<\RatEmbed(\IntEmbed(n))
    \end{aligned}}
    [\FormulaRefAuto{RationalArchimedean}[thm]]%
}

\FormulaThmDeltaK[Das Quotientenmodell erfüllt die rationalen Körperaxiome]{%
  \begin{aligned}[t]
    &\mathbb Q=\QuotientSet{\RatPairCarrier}{\RatPairRel}\ \land{}\\
    &\RationalOrderedField{\mathbb Q}
  \end{aligned}%
}{RationalQuotientModelSatisfiesOrderedField}{%
  \DeltaRow{Mengen}{\mathbb Q}%
  \DeltaRow{Funktionensymbol}{\RatEmbed}%
  \DeltaRow{Relationsoperatoren}{\leq_{\mathbb Q}}%
}

\emph{Beweisskizze.}
Jede Zeile der Zusammenfassung ist entweder eine Definition oder einer der in
diesem Band bewiesenen und dort angegebenen Sätze. Damit erfüllt das aus
positiven Bruchpaaren gebildete Quotientenmodell genau das durch
\FormulaRefAuto{RationalOrderedFieldAxioms}[def] benannte Aussagenpaket.

\end{document}
